\documentclass[pre,twocolumn,superscriptaddress,nofootinbib]{revtex4-2}
\usepackage{graphicx}
\usepackage{amsmath}
\usepackage{booktabs}
\usepackage[hidelinks]{hyperref}

\begin{document}

\title{The fertile window: how the barrier, the release, the room and the relic of a decompactification depend on
the coefficient of the curling cost}

\author{Emily Smith}
\email{emilysmithcreate@gmail.com}
\affiliation{Independent researcher, Charlottesville, Virginia, USA}

\date{DRAFT, 24 September 2026 (night); the window verdict awaits a running test}

\begin{abstract}
In the family of graph energies obtained from combinatorial quantum gravity (CQG) by multiplying its local term
by a coefficient $\lambda$ (only $\lambda = 1$ is CQG; everything here is at $\lambda > 1$, next to CQG and not in
it), a torus with one direction curled is metastable above the flat torus and opens into it by a single front.
We map the ingredients of that change along $\lambda$. Exact: the gap is $4(\lambda - 1)$ per vertex for each
curled direction, in two and in three dimensions; the cheapest single move out of the curled torus costs
$32 - 16\lambda$ and is offered three times a sweep at every size, so the mean waiting time follows an Arrhenius
law with nothing fitted. Measured, pre-registered: the curled torus is metastable at $g = 1.5$ from
$\lambda = 1.05$ to $1.35$ and not at $1.40$; wherever it is metastable the change is a single front with the
two orders side by side, the waits are memoryless, the first exit arrives on time, and about one exit in three
falls back before the front takes hold, at every $\lambda$; the release is exactly $4(\lambda - 1)$; and the
break-up of the single front begins at the edge of the window. Sealed, the reservoir a conversion needs grows
faster than the release, $12$, $36$ and $72$ stores at $\lambda = 1.10$, $1.25$, $1.40$ for 96 vertices, and
the share of decays leaving a relic rises with $\lambda$ while near $\lambda = 1$ none do. In three dimensions,
where the author's rule that all directions curl together can be posed, the fully curled state is metastable
only for $1 < \lambda < 1.2$ (exact, with a derived energy that has not yet passed its reproduction gate). A
larger $\lambda$ gives a larger release, a harder birth and a shorter-lived curled state.
\end{abstract}

\maketitle

\section{Introduction}

The companion paper \cite{paper1} studied one setting, $\lambda = 1.25$, and found that a $4 \times L$ torus with
one direction curled is metastable, opens into the flat torus by nucleation and a single front, and releases
exactly $4(\lambda - 1)$ per vertex. A single setting could be a lucky point. This paper asks how each ingredient
of the change, the barrier, the release, the room the product needs, and the relic left behind, depends on
$\lambda$, from just above the CQG value to the point where the curled torus stops being metastable.

Why a coefficient at all. At $\lambda = 1$ the local term of CQG \cite{T17,KTB19,T25} exactly cancels the reward
for a third square on an edge, so the flat torus, the curled torus and the 4-cube are degenerate: no ordered
arrangement can turn into another with any release. Above 1 the flat torus is the unique ground state and each
curled direction costs a known energy. In the author's framework the extra term is a curling cost, a constant of
the arrangement, and its value is whatever reality has; in a picture in which many such systems exist, the value
we observe would be the one that makes the most productive change. That framework is the author's and is not a
claim of this paper. What this paper measures is the map: at which $\lambda$ the curled state is stuck for now,
how sharply it changes, how much it releases, how much room the release needs, and what it leaves. Every verdict
quoted was pre-registered before its run.

\section{Exact results}

\subsection{The energy, edge by edge}

With $S$ the number of squares and $X = \sum_e (S_e - 2)_+$ the surplus squares on edges carrying more than two,
\begin{equation}
  H = 16(N - S) + 4\lambda X = 4 \sum_e \big[(2 - S_e) + \lambda\,(S_e - 2)_+\big].
  \label{eq:H}
\end{equation}
The flat torus (two squares on every edge) has $H = 0$; a direction curled to a single square puts three squares
on the edges along it, and the ladder is additive: one curled direction costs $4(\lambda - 1)$ per vertex, two
$8(\lambda - 1)$. In three dimensions, with six links per vertex, the same holds with a third rung, so the
fully curled 6-cube sits $12(\lambda - 1)$ per vertex above flat space and each curled direction holds a third
of the release. (The three-dimensional energy is derived by us from the papers' definitions, not quoted, and its
reproduction gate against the published curve has not yet passed; every three-dimensional statement here is exact
arithmetic on that derived energy and nothing more.)

\subsection{The cheapest ways out}

Every switch the chain can propose out of the curled torus was listed and priced. Two families are cheapest:
move A, losing two squares and four surplus squares, at $32 - 16\lambda$, offered $3N$ ways; and move B, at
$64 - 40\lambda$, offered $2N$ ways. A sweep is $2N$ attempts out of $4N^2$ equally likely proposals, so A is
offered three times a sweep and B twice, at every size, and the mean waiting time before the first exit at
coupling $g$ is predicted with nothing fitted:
\begin{equation}
  \tau(\lambda) = \Big[\,3\,e^{-(32 - 16\lambda)/g} + 2\,e^{-(64 - 40\lambda)/g}\Big]^{-1}.
  \label{eq:tau}
\end{equation}
At $g = 1.5$: $\tau = 8330$, $4844$, $2788$, $1570$, $845$, $419$, $183$, $68$ sweeps at
$\lambda = 1.05, 1.10, \dots, 1.40$. Move A is the cheapest way out below $\lambda = 4/3$; above it B takes over.

\subsection{The relic's share, exact}

The common relic is one column of the torus left curled, costing $24\lambda - 16$ above the flat torus. As a share
of the release $4(\lambda - 1)N$ it is $(24\lambda - 16)/[4(\lambda - 1)N]$: 22\% at $\lambda = 1.25$ and 72\% at
$\lambda = 1.05$ for $N = 64$, falling as $1/N$. Below $\lambda = 1 + 8/(4N - 24)$ a column left behind would cost
more than the whole release, so the torus could not end on it downhill.

\subsection{Three dimensions: where the fully curled state is stuck}

The author's rule for the real thing is that the three space directions are intertwined, curl together and open
together. In two dimensions that rule cannot be posed: a gas of fully curled 4-cubes has no wall at any
$\lambda > 1$ (a move inside a knot costs $32 - 32\lambda$, downhill for every $\lambda > 1$), so it falls apart
everywhere at once with no wait. In three dimensions, listing every single move out of each curled torus with the
derived energy: one curled direction is metastable for $1 < \lambda < 2$, two for $1 < \lambda < 1.2$, and all
three (the 6-cube) for $1 < \lambda < 1.2$, with a wall of $96 - 80\lambda$ offered $6.67$ times a sweep. So the
fully curled state can be stuck for now in three dimensions, but only in a narrow window that does not contain
the two-dimensional value of about $1.25$.

\section{The window, measured}

\subsection{Protocol}

The curled $4 \times L$ torus at $g = 1.5$, $N = 64, 96, 144, 192$, followed in blocks of five sweeps until 98\%
converted or 100,000 sweeps; the local-dimension census and the pieces of converted vertices read at 25, 50 and
75\% conversion; the release settled over windows of 600 sweeps; the final graph saved. Three pre-registered
passes: T8 with 30 decays per cell at $\lambda = 1.05$ to $1.45$; T23 with 120 decays per cell at 1.05 to 1.35,
fresh seeds and the memoryless check sized to the sample; and T24, running as this draft is written, with fresh
seeds again and the energy check read exactly from each decay's saved wiring.

\subsection{Metastability and its edge}

The torus is metastable (more than half the decays still a torus when the 200-sweep resting stretch ends) at
every size from $\lambda = 1.05$ to $1.35$, and not at $1.40$ or $1.45$; the edge lies between 1.35 and 1.40 at
$g = 1.5$. That is where the predicted wait, Eq.~\eqref{eq:tau}, falls through the length of the resting stretch,
not where the barrier vanishes: the barriers at the edge are still 8 to 10 units.

\subsection{Sharpness across the window}

At every metastable cell the two orders sit side by side (95 to 100\% of vertices at $d \in \{1, 2\}$ at half
conversion) and the converted region is one piece (76 to 100\% of it, T8; 81 to 100\%, T23). The waits after the
resting stretch are memoryless: in T23 the coefficient of variation lies inside the central 99.9\% band of an
exponential sample of the same size in 20 of 20 window cells from 1.05 to 1.25 and in three of four sizes at 1.30.
The release is exactly $4(\lambda - 1)$ wherever the flat torus is reached. T8 and T23 each came out INCONCLUSIVE
by the letter of their own rules, for stated reasons: T8's memoryless band was too tight for thirty decays, and
T23's energy check, which compared a window average with the exact energy of the final state, tripped on the
sheet's own thermal excitations at $\lambda = 1.30$ in a few decays per hundred. T24 fixes the second before its
run, as T23 fixed the first; its verdict will be reported here whichever way it falls.

\subsection{The first exit is on time, and one exit in three falls back}

Counted move by move, with every accepted move seen (T22, forty decays per cell at $N = 64$ and $96$): the mean time
to the first exit agrees with Eq.~\eqref{eq:tau} within one standard error in all six cells at $\lambda = 1.05$,
$1.10$ and $1.25$, and the number of exits per decay is $1.5$ to $1.8$, so the transmission coefficient
$\kappa = 1/E$ is $0.56$ to $0.68$ with no trend in $\lambda$. Recrossing is a general feature of the chain, not a
near-$\lambda = 1$ effect.

\subsection{The edge of the window: break-up begins there}

Pooled over sizes from $\lambda = 1.30$ to $1.35$ (T23): the share of decays ending at the flat torus falls from
$0.71$ to $0.47$, and the share whose converted region is in more than one piece at 25\% conversion rises from
$0.44$ to $0.69$, each by more than two standard errors. The author predicted this before the run. The single
front gives way to several seeds as the wall comes down.

\section{The room the new space needs}

\emph{Pre-registered as T18.} One seed in a sealed $24 \times 4$ torus, ten bath sizes, twenty runs each, at
$\lambda = 1.10$, $1.25$, $1.40$. The smallest bath giving a majority of clean sheets is $12$, $36$ and $72$ stores.
The release grows as $4(\lambda - 1)$; the room needed grows faster ($C^*/[4(\lambda - 1)N] = 0.31$, $0.375$,
$0.47$). The verdict on record is PROPORTIONAL, with the stated ratio. A larger $\lambda$ gives a larger release
and a harder birth.

\section{The relic along the window}

At fixed coupling, near $\lambda = 1$ no decay leaves a relic, and the share that does rises with $\lambda$ (T8,
T23). Sealed, the relic is one column of fixed size at every $\lambda$ tried, with its share of the release given
exactly above. The companion paper on the relic \cite{paper2} measures its number, its position and its lifetime.

\section{Discussion}

The map has one shape. As $\lambda$ rises the release grows linearly, the wall falls linearly, the wait falls
exponentially, the room the product needs grows faster than the release, and the scrap grows. As $\lambda$ falls
toward 1 everything the change produces shrinks to nothing and the curled state lasts longer. There is no setting
at which the release is large and the curled state is long-lived and the product needs little room: the
ingredients trade against one another across the whole window.

In the author's framework this trade-off is what a selection among loops would select within; the model can map
the ingredients, not count offspring, and this paper claims only the map. In three dimensions the window is
narrower and does not contain the two-dimensional value; whether the released energy of one direction can pay the
wall of the next, the energetic room for a cascade, is exact arithmetic that awaits a kernel that has passed its
reproduction gate.

What this cannot show: anything at exactly $\lambda = 1$, where nothing is released; anything at couplings other
than $g = 1.5$; anything with the points treated as interchangeable, where the wait is expected to grow by a
factor of order $N$ while the release and the front are unchanged; and anything about the real universe.

\section*{Acknowledgments}

Code, analysis and text were drafted with Claude (Anthropic) in conversation with the author, who directed and
checked the work; no physicist has reviewed it. The repository holds every configuration, seed, result and
pre-registration: \url{https://github.com/EmilySmithCreate/RecreationalPhysics}.

\begin{thebibliography}{9}
\bibitem{paper1} E. Smith, The curled torus burps: front-driven decompactification between ordered states in a
graph model of emergent geometry, draft (2026), \texttt{docs/papers/curled\_torus/}.
\bibitem{paper2} E. Smith, What the burp leaves behind, draft (2026), \texttt{docs/papers/relic/}.
\bibitem{T17} C. A. Trugenberger, Combinatorial quantum gravity: geometry from random bits, JHEP 09, 045 (2017).
\bibitem{KTB19} C. Kelly, C. A. Trugenberger and F. Biancalana, Self-assembly of geometric space from random
graphs, Class. Quantum Grav. 36, 125012 (2019).
\bibitem{T25} C. A. Trugenberger, Combinatorial quantum gravity and emergent 3D quantum behaviour, review (2025);
arXiv:2512.17676.
\end{thebibliography}

\end{document}
